\chapter{2005年上海卷（春）}
\section{填空题}
\begin{problem}
方程 \( \lg x^{2}-\lg(x+2)=0 \) 的解集是\fillinblank{}.
\end{problem}

\begin{answer}
\(\{-1,2\}\)
\end{answer}

\begin{solution}
由对数的定义域，得 \(x^{2}>0\) 且 \(x+2>0\)，所以 \(x>-2\) 且 \(x\neq0\). 原方程等价于 \(\lg\dfrac{x^{2}}{x+2}=0\)，从而 \(\dfrac{x^{2}}{x+2}=1\)，即 \(x^{2}-x-2=0\). 因此 \(x=2\) 或 \(x=-1\)，两根都满足定义域条件，故解集为 \(\{-1,2\}\).
\end{solution}

\begin{problem}
\(\displaystyle\lim_{n\to \infty}\frac{n + 2}{1 + 2 + \cdots + n} =\)\fillinblank{}.
\end{problem}

\begin{answer}
0
\end{answer}

\begin{solution}
因为 \(1+2+\cdots+n=\dfrac{n(n+1)}{2}\)，所以
\(\displaystyle\frac{n+2}{1+2+\cdots+n}=\frac{2(n+2)}{n(n+1)}=\frac{2(1+2/n)}{n+1}\).
当 \(n\to\infty\) 时，分子趋于 \(2\)，分母趋于 \(+\infty\)，故极限为 \(0\).
\end{solution}

\begin{problem}
若 \(\cos \alpha = \frac{3}{5}\)，且 \(\alpha \in \left(0, \frac{\pi}{2}\right)\)，则 \(\tan \frac{\alpha}{2} =\)\fillinblank{}.
\end{problem}

\begin{answer}
\(\dfrac{1}{2}\)
\end{answer}

\begin{solution}
由于 \(\alpha\in\left(0,\dfrac{\pi}{2}\right)\)，\(\sin\alpha>0\)，所以
\(\sin\alpha=\sqrt{1-\cos^{2}\alpha}=\dfrac45\). 又由半角公式，
\(\displaystyle\tan\frac\alpha2=\frac{\sin\alpha}{1+\cos\alpha}=\frac{4/5}{1+3/5}=\frac12\).
\end{solution}

\begin{problem}
函数 \( f(x) = -x^{2} \)（\(x \in (-\infty, -2]\)）的反函数 \( f^{-1}(x) = \)\fillinblank{}.
\end{problem}

\begin{answer}
\(f^{-1}(x)=-\sqrt{-x}\)（\(x\leqslant-4\)）
\end{answer}

\begin{solution}
由 \(y=-x^{2}\) 且 \(x\leqslant-2\)，可知值域为 \(y\leqslant-4\). 交换 \(x\)，\(y\) 得 \(x=-y^{2}\)，即 \(y^{2}=-x\). 因为原函数的自变量满足 \(y\leqslant-2\)，所以应取负平方根，得到 \(y=-\sqrt{-x}\)，且 \(x\leqslant-4\). 因此 \(f^{-1}(x)=-\sqrt{-x}\)（\(x\leqslant-4\)）.
\end{solution}

\begin{problem}
在 \(\triangle ABC\) 中，若 \(\angle C = 90^{\circ}\)，\(AC = BC = 4\)，则 \(\overrightarrow{BA} \cdot \overrightarrow{BC} =\)\fillinblank{}.
\end{problem}

\begin{answer}
16
\end{answer}

\begin{solution}
取直角坐标系使 \(C=(0,0)\)，\(A=(4,0)\)，\(B=(0,4)\). 则
\(\overrightarrow{BA}=A-B=(4,-4)\)，\(\overrightarrow{BC}=C-B=(0,-4)\)，从而
\(\overrightarrow{BA}\cdot\overrightarrow{BC}=4\times0+(-4)\times(-4)=16\).
\end{solution}

\begin{problem}
某班共有 \(40\) 名学生，其中只有一对双胞胎，若从中一次随机抽查三位学生的作业，则这对双胞胎的作业同时被抽中的概率是\fillinblank{}. （结果用最简分数表示）
\end{problem}

\begin{answer}
\(\dfrac{1}{260}\)
\end{answer}

\begin{solution}
从 \(40\) 名学生中抽取三人，共有 \(\mathrm C_{40}^{3}\) 种等可能结果. 若双胞胎同时被抽中，第三人可从其余 \(38\) 人中任选，有 \(\mathrm C_{38}^{1}\) 种结果. 因此所求概率为 \(\displaystyle\frac{\mathrm C_{38}^{1}}{\mathrm C_{40}^{3}}=\frac{38}{40\times39\times38/6}=\frac{1}{260}\).
\end{solution}

\begin{problem}
双曲线 \( 9x^{2}-16y^{2}=1 \) 的焦距是\fillinblank{}.
\end{problem}

\begin{answer}
\(\dfrac{5}{6}\)
\end{answer}

\begin{solution}
将方程化为 \(\displaystyle\frac{x^{2}}{1/9}-\frac{y^{2}}{1/16}=1\)，所以
\(a=\dfrac13\)，\(b=\dfrac14\). 双曲线的焦半距参数满足
\(c^{2}=a^{2}+b^{2}\)，因此
\(\displaystyle c=\sqrt{\frac19+\frac1{16}}=\sqrt{\frac{25}{144}}=\frac5{12}\).
焦距是两焦点间的距离 \(2c\)，故焦距为 \(\dfrac56\).
\end{solution}

\begin{problem}
若 \((x+2)^{n}=x^{n}+\cdots+ax^{3}+bx^{2}+cx+2^{n}\)（\(n\in\N\)，\(n\geqslant3\)），且 \(a:b=3:2\)，则 \(n=\)\fillinblank{}.
\end{problem}

\begin{answer}
11
\end{answer}

\begin{solution}
二项式展开式中，\(x^{3}\) 的系数为 \(a=\mathrm C_{n}^{3}2^{n-3}\)，\(x^{2}\) 的系数为 \(b=\mathrm C_{n}^{2}2^{n-2}\). 因而 \(\displaystyle\frac{a}{b}=\frac{\mathrm C_{n}^{3}2^{n-3}}{\mathrm C_{n}^{2}2^{n-2}}=\frac{n-2}{6}\). 由 \(a:b=3:2\)，得 \(\dfrac{n-2}{6}=\dfrac{3}{2}\)，所以 \(n=11\).
\end{solution}

\begin{problem}
设数列 \(\{a_n\}\) 的前 \(n\) 项和为 \(S_{n}\)（\(n \in \N\)）.关于数列 \(\{a_n\}\) 有下列三个命题：

\begin{circlelist}
\item 若 \(\{a_{n}\}\) 既是等差数列又是等比数列，则 \(a_{n} = a_{n + 1}\)（\(n \in \N\)）；
\item 若 \(S_{n} = an^{2} + bn\)（\(a\)，\(b \in \R\)），则 \(\{a_{n}\}\) 是等差数列；
\item 若 \(S_{n} = 1 - (-1)^{n}\)，则 \(\{a_{n}\}\) 是等比数列.
\end{circlelist}

这些命题中，真命题的序号是\fillinblank{}.
\end{problem}

\begin{answer}
\(\circled{1}\circled{2}\circled{3}\)
\end{answer}

\begin{solution}
设首项为 \(a_1\)，公差为 \(d\). 若 \(a_1=0\)，因为它同时是等比数列，\(a_2=a_1q=0\)，从而各项均为 \(0\)，命题\(\circled{1}\)成立. 若 \(a_1\neq0\)，等差数列和等比数列的前三项分别满足
\(a_2=a_1+d\)，\(a_3=a_1+2d\) 以及 \(a_2^2=a_1a_3\). 因此
\((a_1+d)^2=a_1(a_1+2d)\)，化简得 \(d^2=0\)，即 \(d=0\). 此时各项相等，故命题\(\circled{1}\)正确.

令 \(a_0=0\). 由 \(S_n=an^{2}+bn\)，对 \(n\in\N\) 有
\[
a_n=S_n-S_{n-1}=an^2+bn-\bigl(a(n-1)^2+b(n-1)\bigr)=a(2n-1)+b
\]
这是关于 \(n\) 的一次式，相邻两项之差恒为 \(2a\)，所以 \(\{a_n\}\) 是等差数列，命题\(\circled{2}\)正确.

由 \(S_n=1-(-1)^n\)，并令 \(S_0=0\)，得
\(a_n=S_n-S_{n-1}=2(-1)^{n+1}\). 于是
\(\displaystyle\frac{a_{n+1}}{a_n}=-1\)，各项均非零，故 \(\{a_n\}\) 是等比数列，命题\(\circled{3}\)正确. 因此真命题的序号是 \(\circled{1}\)，\(\circled{2}\)，\(\circled{3}\).
\end{solution}

\begin{problem}
若集合 \(A = \{x \mid 3\cos 2\pi x = 3^x, x \in \R\}\)，集合 \(B = \{y \mid y^2 = 1, y \in \R\}\)，则 \(A \cap B =\fillinblank{}\).
\end{problem}

\begin{answer}
\(\{1\}\)
\end{answer}

\begin{solution}
由 \(y^{2}=1\) 且 \(y\in\R\)，得 \(B=\{-1,1\}\). 当 \(x=-1\) 时，\(3\cos(-2\pi)=3\)，而 \(3^{-1}=\dfrac{1}{3}\)，不满足方程；当 \(x=1\) 时，\(3\cos(2\pi)=3=3^{1}\)，满足方程. 所以 \(A\cap B=\{1\}\).
\end{solution}

\begin{problem}
函数 \( y = \sin x + \arcsin x \) 的值域是\fillinblank{}.
\end{problem}

\begin{answer}
\(\left[-\dfrac{\pi}{2}-\sin1,\dfrac{\pi}{2}+\sin1\right]\)
\end{answer}

\begin{solution}
函数的定义域由 \(\arcsin x\) 确定为 \([-1,1]\). 在 \([-1,1]\) 上，\(\sin x\) 严格递增，\(\arcsin x\) 也严格递增，所以两者之和
\(y=\sin x+\arcsin x\) 在该区间上严格递增. 因此
\(y(-1)=-\sin1-\dfrac\pi2\) 为最小值，
\(y(1)=\sin1+\dfrac\pi2\) 为最大值，值域为
\(\left[-\dfrac\pi2-\sin1,\dfrac\pi2+\sin1\right]\).
\end{solution}

\begin{problem}
已知函数 \( f(x) = 2^x + \log_2 x \)，数列 \(\{a_n\}\) 的通项公式是 \( a_n = 0.1n \)（\(n \in \N\)），当 \( |f(a_n) - 2005| \) 取得最小值时，\( n = \)\fillinblank{}.
\end{problem}

\begin{answer}
110
\end{answer}

\begin{solution}
在 \(x>0\) 上，
\(\displaystyle f'(x)=2^x\ln2+\frac1{x\ln2}>0\)，所以 \(f\) 严格递增.
由于 \(a_n=0.1n\)，只需比较 \(a_{110}=11\) 与相邻的 \(a_{109}=10.9\). 计算得
\(f(10.9)=2^{10.9}+\log_2(10.9)\approx1914.30\)，
\(f(11)=2^{11}+\log_2 11\approx2051.46\)，于是
\[
2005-f(10.9)\approx90.70>46.46\approx f(11)-2005
\]
因此 \(f(11)\) 比 \(f(10.9)\) 更接近 \(2005\). 对 \(n\leqslant109\)，由单调性有
\(f(a_n)\leqslant f(10.9)<2005\)，其距离不小于 \(2005-f(10.9)\)；对 \(n\geqslant111\)，有
\(f(a_n)\geqslant f(11.1)>f(11)>2005\)，其距离大于 \(f(11)-2005\). 所以最小值在 \(a_{110}=11\) 时取得，故 \(n=110\).
\end{solution}

\section{单选题}
\begin{problem}
已知直线 \(l\)，\(m\)，\(n\) 及平面 \(\alpha\)，下列命题中的假命题是（\quad）

\choices
    {若 \(l \parallel m\)，\(m \parallel n\)，则 \(l \parallel n\)}
    {若 \(l \perp \alpha\)，\(n \parallel \alpha\)，则 \(l \perp n\)}
    {若 \(l \perp m\)，\(m \parallel n\)，则 \(l \perp n\)}
    {若 \(l \parallel \alpha\)，\(n \parallel \alpha\)，则 \(l \parallel n\)}
\end{problem}

\begin{answer}
D
\end{answer}

\begin{solution}
直线的平行关系具有传递性，所以命题 A 正确；若 \(l\perp\alpha\)，而 \(n\parallel\alpha\)，则 \(n\) 的方向与平面 \(\alpha\) 平行，故 \(l\perp n\)，命题 B 正确. 直线 \(m\) 与 \(n\) 平行时，\(l\perp m\) 可推出 \(l\perp n\)，命题 C 正确. 但两条分别平行于同一平面的直线不一定互相平行，例如它们可以有不同方向，故命题 D 错误. 因此假命题为 D.
\end{solution}

\begin{problem}
在 \(\triangle ABC\) 中，若 \(\frac{a}{\cos A} = \frac{b}{\cos B} = \frac{c}{\cos C}\)，则 \(\triangle ABC\) 是（\quad）

\choices
    {直角三角形}
    {等边三角形}
    {钝角三角形}
    {等腰直角三角形}
\end{problem}

\begin{answer}
B
\end{answer}

\begin{solution}
三角形的角不能为直角，否则相应的余弦为零，题设比值无意义. 若有一个角为钝角，则该角的余弦为负，而另外两个角为锐角、余弦为正；由于 \(a\)，\(b\)，\(c>0\)，三个比值不可能相等. 因而 \(A\)，\(B\)，\(C\) 均为锐角.
由正弦定理，存在外接圆半径 \(R\)，使得
\(a=2R\sin A\)，\(b=2R\sin B\)，\(c=2R\sin C\). 题设等式遂等价于
\(\tan A=\tan B=\tan C\). 正切函数在 \(\left(0,\dfrac\pi2\right)\) 上严格递增，所以
\(A=B=C\). 又 \(A+B+C=\pi\)，故
\(A=B=C=\dfrac\pi3\)，三角形为等边三角形，选 B.
\end{solution}

\begin{problem}
若 \(a\)，\(b\)，\(c\) 是常数，则“\(a > 0\) 且 \(b^2 - 4ac < 0\)”是对“任意 \(x \in \R\)，有 \(ax^2 + bx + c > 0\)”的（\quad）

\choices
    {充分不必要条件}
    {必要不充分条件}
    {充要条件}
    {既不充分也不必要条件}
\end{problem}

\begin{answer}
A
\end{answer}

\begin{solution}
若 \(a>0\) 且 \(\Delta=b^{2}-4ac<0\)，二次函数的图象开口向上且没有与 \(x\) 轴的交点，因此对任意 \(x\in\R\)，都有
\(ax^{2}+bx+c>0\)，题设条件是充分条件.
但它不是必要条件. 例如取 \(a=0\)，\(b=0\)，\(c=1\)，则对任意 \(x\in\R\) 有
\(ax^{2}+bx+c=1>0\)，而 \(a>0\) 不成立. 因此该条件是充分不必要条件，选 A.
\end{solution}

\begin{problem}
设函数 \( f(x) \) 的定义域为 \( \R \)，有下列三个命题：

\begin{circlelist}
\item 若存在常数 \(M\)，使得对任意 \(x \in \R\)，有 \(f(x) \leqslant M\)，则 \(M\) 是函数 \(f(x)\) 的最大值；
\item 若存在 \(x_0 \in \R\)，使得对任意 \(x \in \R\)，且 \(x \neq x_0\)，有 \(f(x) < f(x_0)\)，则 \(f(x_0)\) 是函数 \(f(x)\) 的最大值；
\item 若存在 \(x_0 \in \R\)，使得对任意 \(x \in \R\)，有 \(f(x) \leqslant f(x_0)\)，则 \(f(x_0)\) 是函数 \(f(x)\) 的最小值.
\end{circlelist}

这些命题中，真命题的个数是（\quad）

\choices
    {\(0\) 个}
    {\(1\) 个}
    {\(2\) 个}
    {\(3\) 个}
\end{problem}

\begin{answer}
B
\end{answer}

\begin{solution}
第一个命题不正确. 例如取 \(f(x)=0\)，常数 \(M=1\) 满足对任意 \(x\in\R\) 有 \(f(x)\leqslant M\)，但函数的最大值是 \(0\)，不是 \(1\).

第二个命题正确. 当 \(x=x_0\) 时有 \(f(x)=f(x_0)\)；当 \(x\neq x_0\) 时题设给出 \(f(x)<f(x_0)\). 因而对所有 \(x\in\R\) 都有 \(f(x)\leqslant f(x_0)\)，且等号在 \(x=x_0\) 处取得，所以 \(f(x_0)\) 是最大值.

第三个命题不正确. 题设不等式恰好说明 \(f(x_0)\) 是最大值，而不能说明它是最小值. 例如取 \(f(x)=-(x-1)^2\)，\(x_0=1\)，则对任意 \(x\in\R\) 有 \(f(x)\leqslant f(1)=0\)，但函数没有最小值，故 \(f(1)\) 不是最小值. 因此只有一个命题为真，选 B.
\end{solution}

\section{解答题}
\begin{problem}
已知 \(z\) 是复数，\(z + 2\i\)，\(\frac{z}{2 - \i}\) 均为实数（\(\i\) 为虚数单位），且复数 \((z + a\i)^2\) 在复平面上对应的点在第一象限，求实数 \(a\) 的取值范围.
\end{problem}

\begin{solution}
设 \(z=x+y\i\)，其中 \(x\)，\(y\in\R\). 因为 \(z+2\i\) 为实数，其虚部为零，所以 \(y+2=0\)，即 \(y=-2\)，从而 \(z=x-2\i\).

又
\[
\frac{z}{2-\i}=\frac{(x-2\i)(2+\i)}{(2-\i)(2+\i)}
=\frac{2x+2+(x-4)\i}{5}
=\frac{2x+2}{5}+\frac{x-4}{5}\i
\]
该复数为实数，当且仅当虚部为零，即 \(x-4=0\). 因此 \(z=4-2\i\).

于是
\[
(z+a\i)^2=\left(4+(a-2)\i\right)^2
=\left(16-(a-2)^2\right)+8(a-2)\i
\]
复平面上对应的点在第一象限，当且仅当实部、虚部均为正，即
\(16-(a-2)^2>0\) 且 \(8(a-2)>0\). 前一个不等式等价于
\(-2<a<6\)，后一个不等式等价于 \(a>2\). 联立得
\(2<a<6\).
\end{solution}

\begin{problem}
已知正三棱锥 \(P - ABC\) 的体积为 \(72\sqrt{3}\)，侧面与底面所成的二面角的大小为 \(60^{\circ}\).

\begin{enumerate}
\item 证明： \(PA \perp BC\)；
\item 求底面中心 \(O\) 到侧面的距离.
\end{enumerate}
\end{problem}

\begin{solution}
设正三角形底面 \(ABC\) 的边长为 \(s\)，中心为 \(O\)，\(M\) 为 \(AB\) 的中点. 正三棱锥的高为 \(PO\)，且 \(PO\perp\) 平面 \(ABC\). 由于 \(P\) 在底面中心 \(O\) 的正上方，侧面三角形 \(PAB\) 为等腰三角形，所以 \(PM\perp AB\)；同时 \(OM\perp AB\). 因而平面 \(POM\perp AB\)，侧面与底面所成的二面角为 \(\angle PMO=60^\circ\).

正三角形的内切半径为
\(\displaystyle OM=\frac{\sqrt3}{6}s\). 在直角三角形 \(POM\) 中，
\(\displaystyle PO=OM\tan60^\circ=\frac{s}{2}\).
底面面积为 \(\displaystyle\frac{\sqrt3}{4}s^2\)，所以由体积条件
\[
72\sqrt3=\frac13\cdot\frac{\sqrt3}{4}s^2\cdot\frac{s}{2}
=\frac{\sqrt3}{24}s^3
\]
故 \(s^3=1728\)，且 \(s>0\)，从而 \(s=12\). 因此
\(PO=6\)，\(OM=2\sqrt3\)，并由勾股定理得
\(PM=\sqrt{PO^2+OM^2}=\sqrt{36+12}=4\sqrt3\).

\begin{enumerate}
\item 在底面正三角形 \(ABC\) 中，中心线 \(AO\) 垂直于边 \(BC\). 又 \(PO\perp\) 平面 \(ABC\)，故 \(PO\perp BC\). 直线 \(AO\) 与 \(PO\) 在平面 \(PAO\) 内相交，且 \(BC\) 同时垂直于这两条相交直线，所以 \(BC\perp\) 平面 \(PAO\). 由于 \(PA\subset\) 平面 \(PAO\)，得到 \(PA\perp BC\).

\item 以侧面 \(PAB\) 为例. 平面 \(POM\) 垂直于 \(AB\)，它与平面 \(PAB\) 的交线为 \(PM\)，故点 \(O\) 到平面 \(PAB\) 的距离等于点 \(O\) 到直线 \(PM\) 的距离. 在直角三角形 \(POM\) 中，\(O\) 到斜边 \(PM\) 的高为
\[
d=\frac{PO\cdot OM}{PM}
=\frac{6\cdot2\sqrt3}{4\sqrt3}=3
\]
由正三棱锥的对称性，\(O\) 到任一侧面的距离都相同，所以所求距离为 \(3\).
\end{enumerate}
\end{solution}

\begin{problem}
已知 \(\tan \alpha\) 是方程 \(x^{2} + 2x\sec \alpha + 1 = 0\) 的两个根中较小的根，求 \(\alpha\) 的值.
\end{problem}

\begin{solution}
因为题设说 \(\tan\alpha\) 是根，所以 \(\tan\alpha\) 有定义，即 \(\cos\alpha\neq0\). 将 \(x=\tan\alpha\) 代入方程，得
\[
\tan^2\alpha+2\tan\alpha\sec\alpha+1
=\sec^2\alpha+2\tan\alpha\sec\alpha
=\sec\alpha\left(\sec\alpha+2\tan\alpha\right)=0
\]
由于 \(\sec\alpha\neq0\)，所以 \(\sec\alpha+2\tan\alpha=0\). 两边同乘 \(\cos\alpha\)，得
\(1+2\sin\alpha=0\)，即 \(\sin\alpha=-\dfrac12\).

因此
\(\displaystyle\alpha=\frac{7\pi}{6}+2k\pi\) 或
\(\displaystyle\alpha=\frac{11\pi}{6}+2k\pi\)，其中 \(k\in\Z\).
当 \(\displaystyle\alpha=\frac{7\pi}{6}+2k\pi\) 时，
\(\displaystyle\tan\alpha=\frac1{\sqrt3}\)，
\(\displaystyle\sec\alpha=-\frac2{\sqrt3}\). 原方程为
\(x^2-\dfrac4{\sqrt3}x+1=0\)，其两根为
\(\displaystyle\frac1{\sqrt3}\) 和 \(\sqrt3\)，所以 \(\tan\alpha\) 是较小根.

当 \(\displaystyle\alpha=\frac{11\pi}{6}+2k\pi\) 时，
\(\displaystyle\tan\alpha=-\frac1{\sqrt3}\)，
\(\displaystyle\sec\alpha=\frac2{\sqrt3}\). 原方程为
\(x^2+\dfrac4{\sqrt3}x+1=0\)，其两根为
\(-\sqrt3\) 和 \(-\dfrac1{\sqrt3}\)，此时 \(\tan\alpha=-\dfrac1{\sqrt3}\) 是较大根，不合题意. 故
\(\displaystyle\alpha=\frac{7\pi}{6}+2k\pi\)，\(k\in\Z\).
\end{solution}

\begin{problem}
某市 \(2004\) 年底有住房面积 \(1200\text{ 万平方米}\)，计划从 \(2005\) 年起，每年拆除 \(20\text{ 万平方米}\) 的旧住房，假定该市每年新建住房面积是上年年底住房面积的 \(5\%\).

\begin{enumerate}
\item 分别求 \(2005\) 年底和 \(2006\) 年底的住房面积；
\item 求 \(2024\) 年底的住房面积. （计算结果以万平方米为单位，且精确到 \(0.01\)）
\end{enumerate}
\end{problem}

\begin{solution}
设 \(A_n\) 表示第 \(2004+n\) 年年底的住房面积，单位为万平方米. 由题意，\(A_0=1200\). 从第 \(2005\) 年起，第 \(n\) 年新建面积为上一年年底面积的 \(5\%\)，同时拆除 \(20\) 万平方米，所以
\(A_n=A_{n-1}+0.05A_{n-1}-20=1.05A_{n-1}-20\).

\begin{enumerate}
\item 由递推式得 \(A_1=1.05\times1200-20=1240\)，\(A_2=1.05\times1240-20=1282\). 因此 2005 年底和 2006 年底的住房面积分别为 \(1240\) 万平方米和 \(1282\) 万平方米.

\item 令 \(B_n=A_n-400\). 由递推式得
\(B_n=1.05A_{n-1}-420=1.05(A_{n-1}-400)=1.05B_{n-1}\)，且 \(B_0=800\). 因此
\(B_n=800\times1.05^n\)，即
\(A_n=400+800\times1.05^n\).
2024 年底对应 \(n=2024-2004=20\)，所以
\(A_{20}=400+800\times1.05^{20}\approx2522.6382\).
精确到 \(0.01\) 得 \(A_{20}\approx2522.64\)，即 2024 年底住房面积为 \(2522.64\) 万平方米.
\end{enumerate}
\end{solution}

\begin{problem}
已知函数 \( f(x) = x + \frac{a}{x} \) 的定义域为 \((0, +\infty)\)，且 \( f(2) = 2 + \frac{\sqrt{2}}{2} \). 设点 \(P\) 是函数图象上的任意一点，过点 \(P\) 分别作直线 \( y = x \) 和 \( y \) 轴的垂线，垂足分别为 \(M\)，\(N\).

\begin{enumerate}
\item 求 \(a\) 的值；
\item 问： \( |PM| \cdot |PN| \) 是否为定值？若是，则求出该定值；若不是，则说明理由；
\item 设 \(O\) 为坐标原点，求四边形 \(OMPN\) 面积的最小值.
\end{enumerate}

\bitmapfigure[max width=0.42\linewidth]{img/2005/shanghai_spring/q21_fig1.png}
\end{problem}

\begin{solution}
\begin{enumerate}
\item 由 \(f(2)=2+\dfrac{a}{2}=2+\dfrac{\sqrt2}{2}\)，得 \(a=\sqrt2\). 设 \(P=(x,y)\)，则 \(x>0\)，且
\(y=x+\dfrac{\sqrt2}{x}\).

\item 点 \(N\) 是点 \(P\) 在 \(y\) 轴上的垂足，所以 \(N=(0,y)\)，从而 \(PN=x\). 点 \(M\) 是点 \(P\) 到直线 \(y=x\) 的垂足. 由 \(x>0\) 和 \(a=\sqrt2>0\)，有 \(y-x=\dfrac{\sqrt2}{x}>0\)，故点到直线 \(y=x\) 的距离为
\(\displaystyle PM=\frac{y-x}{\sqrt2}=\frac1x\).
于是
\(\displaystyle |PM|\cdot|PN|=\frac1x\cdot x=1\)，与 \(P\) 的位置无关，所以是定值 \(1\).

\item 记 \(d=y-x=\dfrac{\sqrt2}{x}>0\). 由于 \(PM\perp y=x\)，垂足为
\(\displaystyle M=\left(x+\frac d2,x+\frac d2\right)\)，而 \(N=(0,x+d)\).
四边形 \(OMPN\) 的对角线 \(OP\) 将其分成三角形 \(OMP\) 和 \(OPN\). 因此
\[
S_{OMPN}
=\frac12\left(x+\frac d2\right)d+\frac12x(x+d)
=\frac{x^2}{2}+xd+\frac{d^2}{4}
\]
代入 \(\displaystyle d=\frac{\sqrt2}{x}\)，得
\[
S_{OMPN}
=\frac{x^2}{2}+\sqrt2+\frac1{2x^2}
=\sqrt2+\frac12\left(x^2+\frac1{x^2}\right)
\]
因为 \(x>0\)，由基本不等式
\(\displaystyle x^2+\frac1{x^2}\geqslant2\)，等号当且仅当 \(x^2=1\)，结合 \(x>0\) 得 \(x=1\). 所以四边形 \(OMPN\) 面积的最小值为 \(1+\sqrt2\).
\end{enumerate}
\end{solution}

\begin{problem}
\begin{enumerate}
\item 求右焦点坐标是 \((2,0)\)，且经过点 \((-2, -\sqrt{2})\) 的椭圆的标准方程；
\item 已知椭圆 \(C\) 的方程是 \(\frac{x^2}{a^2} + \frac{y^2}{b^2} = 1\)（\(a > b > 0\)）.设斜率为 \(k\) 的直线 \(l\) 交椭圆 \(C\) 于 \(A\)，\(B\) 两点，\(AB\) 的中点为 \(M\).证明：当直线 \(l\) 平行移动时，动点 \(M\) 在一条过原点的定直线上；
\item 利用上一问所揭示的椭圆几何性质，用作图方法找出下面给定椭圆的中心，简要写出作图步骤，并在图中标出椭圆的中心.
\end{enumerate}

\FigureLayoutDeclare{img/2005/shanghai_spring/q22_fig1.png}{5.5cm}{61bp 157bp 106bp 50bp}
\bitmapfigure[float=inline,max width=0.4\linewidth]{img/2005/shanghai_spring/q22_fig1.png}
\end{problem}

\begin{solution}
\begin{enumerate}
\item 右焦点为 \((2,0)\)，所以椭圆焦半距 \(c=2\)，且长轴在 \(x\) 轴上. 设其标准方程为
\(\displaystyle\frac{x^2}{a^2}+\frac{y^2}{b^2}=1\)，其中 \(a > b > 0\). 由
\(c^2=a^2-b^2\) 得 \(a^2-b^2=4\).
又因为 \((-2,-\sqrt2)\) 在椭圆上，
\(\displaystyle\frac4{a^2}+\frac2{b^2}=1\). 令 \(A=a^2\)，则
\(b^2=A-4\)，且 \(A>4\). 代入得
\(\frac4A+\frac2{A-4}=1\)，\(\Longrightarrow\)，\(4(A-4)+2A=A(A-4)\)，\(\Longrightarrow\)，\(A^2-10A+16=0\).
故 \(A=2\) 或 \(A=8\)，由 \(A>4\) 舍去 \(A=2\)，得到
\(a^2=8\)，\(b^2=4\). 所求椭圆方程为
\(\displaystyle\frac{x^2}{8}+\frac{y^2}{4}=1\).

\item 设直线 \(l\) 的方程为 \(y=kx+t\)，其中斜率 \(k\) 固定，截距 \(t\) 随平行移动而变化. 代入椭圆方程并整理，得
\[
(b^2+a^2k^2)x^2+2a^2ktx+a^2(t^2-b^2)=0
\]
设交点 \(A\)，\(B\) 的横坐标分别为 \(x_1\)，\(x_2\)，则由韦达定理，
\[
x_M=\frac{x_1+x_2}{2}
=-\frac{a^2kt}{b^2+a^2k^2}
\]
又因为 \(M\) 在直线 \(l\) 上，
\[
y_M=kx_M+t
=\frac{b^2t}{b^2+a^2k^2}
\]
消去 \(t\)，可得 \(a^2k y_M+b^2x_M=0\). 这是一条过原点的直线，且方程中不含平行移动参数 \(t\)，所以当 \(l\) 平行移动时，弦 \(AB\) 的中点 \(M\) 始终在这条定直线上. 当 \(k=0\) 时，方程变为 \(x_M=0\)，仍是过原点的定直线.

\item 在给定椭圆内任选一个方向，作两条互不重合且与椭圆相交的平行弦，分别作出两条弦的中点，并连接这两个中点. 由上一问，所得直线必经过椭圆中心. 再取一个不同的方向，重复上述作图，得到另一条经过椭圆中心的直线. 两条中点连线不平行，故其交点唯一，该交点就是椭圆中心，在图中将它标为 \(O\).
\end{enumerate}
\end{solution}
